Para instrumentar el banco de pruebas se requiere un enlace digital rápido y simple entre el microcontrolador y dos periféricos clave: el sensor magnético TMAG5170, que entrega el ángulo del imán solidario al eje del motor, y el controlador Ethernet ENC28J60, que transmite los datos al computador. La Interfaz Periférica Serial (SPI) es un protocolo síncrono de baja sobrecarga, adecuado para distancias cortas y muy usado en sistemas embebidos por su sencillez y altas tasas de reloj \autocite{Pini2019_SPI_DigiKey}.

\section{Fundamentos de SPI}
Un bus SPI básico emplea hasta cuatro señales: SCK, MOSI, MISO y una línea CS por dispositivo. El maestro genera el reloj y habilita un esclavo a la vez llevando su CS al nivel activo; los demás permanecen en alta impedancia sobre MISO. SPI es dúplex completo: en cada flanco del reloj puede enviarse un bit por MOSI mientras el esclavo devuelve un bit por MISO. El maestro configura la polaridad y la fase del reloj (modos 0–3); todos los dispositivos del bus deben compartir el mismo modo. Las velocidades típicas alcanzan varios Mhz en placas con pistas cortas y buen plano de masa \cite{Pini2019_SPI_DigiKey}.

\section{Aplicación en este proyecto}
El TMAG5170 y el ENC28J60 comparten SCK, MOSI y MISO; cada uno dispone de su CS dedicada. Se selecciona un único modo de cronometraje y una frecuencia común definida por el dispositivo más lento. El diseño de PCB considera trazas cortas y retorno de masa directo bajo el bus, separación del lazo de potencia del motor y resistencias serie cercanas al maestro cuando sea necesario para amortiguar flancos. Con esta topología se logra una comunicación simple y confiable para adquisición y envío de datos, siguiendo las pautas generales del protocolo descritas por DigiKey \cite{Pini2019_SPI_DigiKey}.
